% Research fragment: independently derived 2026-10-08.
% Theorem environments and amsmath/amssymb are supplied by the parent article.
% Citations are supplied in sources.bib in this directory.

\section{Four-ended rational tangles have exponentially many decision profiles}
\label{sec:rational-continuation-rank}

The following construction separates three different size measures: the
number of exact continuation-equivalence classes, the dimension of a
linear representation of the final Boolean answer, and the number of bits
needed for an arithmetic representation.  The first two are exponential
even at a fixed four-point boundary.  The third is linear on the same
family, and the resulting decision procedure has polynomial bit complexity.

\subsection{Rational tangle conventions and the gluing test}

Write $T(p/q)$ for a rational two-string tangle with reduced fraction $p/q$,
and write $N(T)$ for numerator closure.  All four boundary points are
marked.  A minus sign on a tangle denotes its mirror image.  We use
\[
 F(T+[k])=F(T)+k,
 \qquad F(T^{\rm rot})=-\frac1{F(T)}.
\]
Horizontal addition of two crossings and vertical addition of two crossings
therefore act on a primitive fraction column by
\begin{equation}
 A\binom pq=\binom{p+2q}{q},\qquad
 B\binom pq=\binom p{2p+q},\qquad
 A=\begin{pmatrix}1&2\\0&1\end{pmatrix},\quad
 B=\begin{pmatrix}1&0\\2&1\end{pmatrix}.
 \label{eq:rational-two-twist-generators}
\end{equation}
Here the vertical operation is
$T\mapsto T*1/[2]=1/([2]+1/T)$.
Each operation adds two crossings and preserves the pairing of the four
marked endpoints.  These fraction and tangle operations are reviewed in
\cite[Section~2]{KauffmanLambropoulouHard}.

We use the following classical, exact rational-tangle gluing criterion:
\begin{equation}
 N\bigl(T(p/q)-T(r/s)\bigr)\text{ is the unknot}
 \quad\Longleftrightarrow\quad |ps-qr|=1.
 \label{eq:rational-unknot-gluing}
\end{equation}
This is \cite[Theorem~5]{KauffmanLambropoulouHard}, with the second fraction
negated.  The two input fraction pairs must be primitive separately.
The integer $ps-qr$ is computed before any cancellation in the sum of the
fractions.  For example, $1/2+1/2=1$ as a rational number, whereas the
numerator closure of $T(1/2)+T(1/2)$ is not an unknot: the integer in the
sum version of the criterion is $1\cdot2+2\cdot1=4$.
This distinction is necessary because a sum of rational tangles need not
itself be a rational tangle.

\subsection{A family of disjoint positive cones}

Let $\mathcal C=\mathbb R_{>0}^{2}$.  For a word
$w=w_1\cdots w_m\in\{A,B\}^m$, let
\[
 G_w=w_1\cdots w_m\in\operatorname{SL}_2(\mathbb Z),
 \qquad \mathcal C_w=G_w\mathcal C.
\]
The product convention is literal matrix multiplication.  To construct
the corresponding tangle from its base tangle, apply the letters of $w$
from right to left.

\begin{lemma}[Disjoint cones]
\label{lem:rational-disjoint-cones}
If $u,w\in\{A,B\}^m$ are distinct, then
$\mathcal C_u\cap\mathcal C_w=\varnothing$.
Moreover, if $z\in\mathcal C_u$ and
$y=G_w^{-1}z\in\mathbb Z^2$, then at least one coordinate of $y$ is
nonpositive and at least one coordinate is strictly positive.
\end{lemma}

\begin{proof}
Every generator preserves $\mathcal C$, and
\[
 A\mathcal C=\{(x,y):y>0,\ x>2y\},\qquad
 B\mathcal C=\{(x,y):x>0,\ y>2x\}.
\]
These two sets are disjoint.  Remove the longest common prefix of $u,w$;
its matrix is invertible.  The remaining two words start with different
generators, and their tail matrices preserve $\mathcal C$.  Their cone
images are therefore contained in the two disjoint sets above.  Applying
the common prefix back proves disjointness.

If both coordinates of $y$ were positive, then $z\in\mathcal C_w$,
contradicting disjointness.  If both were nonpositive, the nonnegative
matrix $G_w$ would send $y$ to a vector with both coordinates nonpositive,
contradicting $z\in\mathcal C_u\subset\mathcal C$.  This proves the
asserted sign pattern.  It also handles a zero coordinate explicitly:
the other coordinate must then be a positive integer.
\end{proof}

\begin{theorem}[An exponential identity submatrix using knots only]
\label{thm:rational-boolean-identity}
For every integer $m\ge0$, there are $2^m$ rational prefix tangles $P_w$
and $2^m$ rational continuation tangles $Q_w$, each with the same four
marked boundary points, such that:
\begin{enumerate}
 \item $P_w$ has a diagram with at most $2m+1$ crossings, and $Q_w$
       has a diagram with at most $2m+3$ crossings;
 \item every closure $K_{u,w}=N(P_u+Q_w)$ is a one-component knot;
 \item $K_{u,w}$ is the unknot if and only if $u=w$.
\end{enumerate}
In particular, the matrix
\[
 H_m(u,w)=\mathbf 1\{N(P_u+Q_w)\text{ is the unknot}\}
\]
is the $2^m\times2^m$ identity matrix and has rank $2^m$ over every field.
Every displayed closed diagram has at most $4m+4$ crossings.
\end{theorem}

\begin{proof}
Define primitive columns
\[
 v_w=\binom{p_w}{q_w}=G_w\binom11,
 \qquad
 x_w=\binom{r_w}{s_w}=G_w\binom32.
\]
All coordinates are positive integers.  Primitivity follows because the
base columns are primitive and $G_w$ is unimodular.  Let
\[
 P_w=T(p_w/q_w),\qquad Q_w=-T(r_w/s_w),
\]
using the diagrams obtained from base tangles $T(1)$ and $T(3/2)$ by the
two-crossing operations in \eqref{eq:rational-two-twist-generators}.
The bases have one and three crossings, respectively; the latter is
$[1,2]=[1]+1/[2]$.  Taking a mirror image preserves the crossing count.
The asserted size bounds follow.

Each generator adds an even number of twists to an adjacent pair of
boundary endpoints, so it preserves the endpoint pairing.  Thus every
$P_w$ has the pairing of $T(1)$, and every $Q_w$ has that of $-T(3/2)$.
Rational tangles have two arcs and no closed components.  Joining these
two fixed endpoint pairings and adding the two numerator caps produces
one circle; equivalently, the base closure $N(T(1)-T(3/2))$ is an unknot
by \eqref{eq:rational-unknot-gluing} and all subsequent two-twist changes
preserve its component count.  Hence every cross closure is a knot.
As an independent parity check, $G_w\equiv I\pmod2$, so
\[
 v_w\equiv(1,1)^t,\qquad x_w\equiv(1,0)^t\pmod2,
 \qquad \det(v_u,x_w)\equiv1\pmod2.
\]

For the diagonal entries, unimodularity gives
\[
 \det(v_w,x_w)=\det\left(\binom11,\binom32\right)=-1.
\]
For $u\ne w$, let $y=(a,b)^t=G_w^{-1}v_u$.  By
Lemma~\ref{lem:rational-disjoint-cones}, $a,b$ are integers with one
nonpositive coordinate and one positive coordinate.  Therefore
\[
 \det(v_u,x_w)=\det\left(\binom ab,\binom32\right)=2a-3b.
\]
If $a\le0<b$, this is at most $-3$; if $b\le0<a$, it is at least $2$.
Consequently its absolute value is at least two.  The parity calculation
in fact makes it at least three.  Criterion
\eqref{eq:rational-unknot-gluing} proves the precise diagonal acceptance
pattern.  The rank assertion follows from the identity matrix.
\end{proof}

\begin{remark}[A direct numerical gap]
The generators above give the stronger off-diagonal bound
$|\det(v_u,x_w)|\ge19$.  Remove the longest common prefix.  If the first
different letters are $A$ on the row side and $B$ on the column side,
write the remaining positive columns as $(a,b)^t$ and $(c,d)^t$.
Here $a,b\ge1$, $c\ge3$, $d\ge2$.  Then
\[
 \det\left(A\binom ab,B\binom cd\right)
 =2ac+ad+3bc+2bd\ge21.
\]
In the opposite order the determinant is
$-(2ac+3ad+bc+2bd)\le-19$.
This gap is useful as a simple exact-arithmetic diagnostic; the cone
argument alone already proves the theorem.
\end{remark}

\subsection{Consequences for continuation summaries}

\begin{corollary}[Finite-state and linear-summary obstructions]
\label{cor:rational-summary-obstructions}
Consider an exact summary $\sigma(P)$ of a four-ended prefix tangle from
which the unknot decision can be recovered for every admitted continuation,
and suppose the admitted continuations include the $Q_w$ above.
On prefixes with at most $2m+1$ crossings, $\sigma$ needs at least $2^m$
distinct values.  Any fixed-length binary representation of these values
needs at least $m$ bits.

If, more specifically, there are vectors $f(P)\in\mathbb F^d$ and
linear functionals $\lambda_Q$ satisfying
\[
 \lambda_Q(f(P))=\mathbf 1\{N(P+Q)\text{ is the unknot}\}
\]
for every such prefix and continuation, then $d\ge2^m$.
The same conclusion holds for a bilinear factorization of the Boolean
connection matrix through a $d$-dimensional vector space.
\end{corollary}

\begin{proof}
The continuation $Q_w$ accepts $P_w$ and rejects every other $P_u$ in the
family, so their summaries must all be distinct.  The bit lower bound is
$\log_2(2^m)=m$.  In the linear case the restricted matrix $H_m$ factors
through a vector space of dimension $d$, so its rank is at most $d$.
Apply Theorem~\ref{thm:rational-boolean-identity}.
\end{proof}

This argument concerns exact representations of \emph{all} continuation
answers.  It does not give a running-time lower bound for unknot recognition.
In particular, $2^m$ different states can be represented by $O(m)$ bits.
The four-ended interface is fixed throughout.  Each tangle is built by
twisting two of four boundary strands, and a collar decomposition gives
constant-size intermediate boundaries.  Any use of this observation for a
particular PD sweep convention should state that convention explicitly.

\begin{theorem}[A rank-two arithmetic representation with a nonlinear readout]
\label{thm:rational-arithmetic-representation}
For rational tangles supplied by explicit two-twist words as above,
the pair of integers $(p,q)$ is an exact tangle summary using $O(m)$ bits.
Two such summaries decide the closure unknot question in $O(m^2)$ bit
operations using schoolbook arithmetic and $O(m)$ working bits.

For the family in Theorem~\ref{thm:rational-boolean-identity}, the signed
gluing-value matrix
\[
 D_m(u,w)=p_u s_w-q_u r_w
\]
has rank at most two over $\mathbb Q$, whereas
$H_m(u,w)=\mathbf1\{|D_m(u,w)|=1\}$ has rank $2^m$.
\end{theorem}

\begin{proof}
Both generators have infinity-operator norm at most three.  Thus
\[
 \|v_w\|_\infty\le3^m,\qquad
 \|x_w\|_\infty\le3^{m+1}.
\]
Their coordinates have $O(m)$ binary digits.  Each generator update is one
doubling and one addition of integers having $O(m)$ bits.  Performing all
$m$ updates uses $O(m^2)$ bit operations and $O(m)$ working bits.
The final two multiplications, subtraction and comparison with $\pm1$
also fit these bounds by schoolbook integer multiplication.

Conway's classification of rational tangles by their fraction
\cite[Theorem~1]{kauffman-lambropoulou-tangles} makes $(p,q)$, normalized up
to simultaneous sign, an exact summary of the rational tangle with marked
boundary.  Efficient evaluation of a gluing is supplied here for rational
continuations by \eqref{eq:rational-unknot-gluing}.

Finally, if $p,q,r,s$ denote the column vectors of the corresponding
coordinates across the finite family, then
\[
 D_m=ps^t-qr^t.
\]
This is a sum of two rank-one matrices, so its rank is at most two.
The Boolean-rank assertion is Theorem~\ref{thm:rational-boolean-identity}.
\end{proof}

The new obstruction therefore identifies a limitation of one specific
compression target: linear factorization of the final exact Boolean
predicate.  Keeping the signed arithmetic gluing value before applying
its nonlinear readout avoids that obstruction on rational tangles.
Extending this arithmetic representation beyond certified rational pieces
requires additional topological structure and remains a research question.
